\section{Linear trace and inverse correction in ten positive integers}
\label{sec:positive}
The three equations in this section have exactly ten independent inputs
\[
 J,F,\alpha,i,T,w,s,n,\eta,\zeta\in\Z_{>0}.
\]
Their complete four-port tuple is the Start-only4127 specialization
of Section~\ref{sec:start4127-positive10}. The arithmetic proof below is
uniform under its authenticated scalar bounds; the whole-native bridge
in that section licenses all four ports together. Strict Horizons
supplies the free quotient and projection-first strategy; Trace Radical
supplies the linear trace; Unsquared Factors and Linear Selectors
supplies the inverse correction at $k$. The retained Round21 uniform signed bound
licenses their composition without assuming $q\mid Y$ early and
without the extra positive shift $T+\alpha$.

\subsection{The three complete polynomial definitions}
Put $h=\beta J$, $q=h+1$, $C=q-F-\alpha$ and $k=\eta+\zeta$.
Choose one of the following fixed charts. The symbols $W_b$, $\lambda$
and $D$ below are polynomial abbreviations, not independent inputs.
\begin{center}\small
\begin{tabular}{@{}lccccc@{}}\toprule
Chart & $W_b$ & $X$ & $Y$ & $\lambda$ & $D$\\\midrule
Free plain & $qw$ & $q^2w$ & $s$ & $q$ & $q^7$\\
Scaled inverse & $w$ & $qw$ & $qs$ & $1$ & $6q^6$\\
Free inverse & $w$ & $qw$ & $s$ & $q$ & $6q^7$\\\bottomrule
\end{tabular}
\end{center}
In every chart define
\begin{align*}
 c&=kY+\eta,& A&=\lambda Y(X+1)+1,&\Delta&=A^2-1,\\
 H&=4A-5,&v&=X\lambda^2Y^2,&E&=1+v(v+1)k^2,\\
 R&=[h(q-F)+\alpha](q^2-1)+(M_C+qM_F)J,\\
 B&=(K_0+W_b)C-F,&U&=ic^2,&a&=\Delta U^2E,\\
 F_0&=1+a,&G&=EF_0,&f&=1+2a,\\
 Q&=4\Delta aF_0,&b&=cT-R,&V&=fb-c,\\
 K&=X(3c-2X)+2,&L&=bV,&M&=hDH.
\end{align*}
In the plain chart put $e_c=k-R-1-q^4s$; in either inverse chart put
$e_c=6(k-R-1-q^4s)-vk(k^2-4)$. Finally set
\begin{align}
 Z&=hD[H(n-LB)+K]-H[DB+he_c],\nonumber\\
 P&=Q(Z^2-M^2GV^2)+(M^2G-Z^2).\label{eq:positive-polynomial}
\end{align}
The equation is $P=0$. Fixed powers are implemented by binary
multiplications. There is no radical or quotient gate. The graph's
register sharing, not the typography of these abbreviations, defines
the measured operation count.

For the proof only, let $t_*=B/h$, $\kappa=K/H$, $z=Z/M$ and
$W_*=1+hL$. In the plain chart let $e=k-R-1$, $\rho=e/q^7$.
In the inverse charts let
\[
 e_*=k-R-1-\frac{vk(k^2-4)}6,
 \qquad \rho=e_*/q^d,
 \quad d=6\text{ or }7.
\]
In all charts put $\sigma=Y/q^3$. Direct expansion gives
\begin{equation}\label{eq:positive-carrier}
 P=(Q-1)Z^2-M^2G(QV^2-1),\qquad
 z=n+\kappa+\sigma-W_*t_*-\rho.
\end{equation}
These proof quotients are rational until integrality is established.
The carrier denominator is nonzero for every positive input.

\subsection{Root extraction before any transport sign}
Positive inputs give $q\ge32$, $X\ge q$, $Y\ge1$, $k\ge2$,
$c\ge3$, $A>q^2$, $\Delta>1$, $U\ge c^2$ and $E>1$.
The exact identities are
\begin{equation}\label{eq:positive-trace}
 f^2=1+4aF_0,\quad Q=\Delta(f^2-1)=4\Delta^2U^2G,
 \quad\gcd(E,F_0)=1.
\end{equation}
Also $f>2c$ and $Q>c^2$, hence $V\ne0$ for every integral $b$.
At a zero, valuations in \eqref{eq:positive-carrier} imply that
$G=g^2$ for a positive integer $g$: every prime of $G$ divides neither
$Q-1$ nor $QV^2-1$, so its valuation in the rational square $z^2$ is
even. Coprimality then gives $E=\tau^2$, $F_0=f_0^2$ and
$g=\tau f_0$ with positive integer roots.

The inherited adjacent-coefficient descent makes $y=|z|/g$ integral.
For completeness, a reduced denominator $m$ satisfies $m^2\mid Q-1$.
For $u=|V|>1$, the norm-preserving transformation
\[
 u'=(2Q-1)u-2(Q-1)y,\qquad
 y'=(2Q-1)y-2Qu
\]
has a nonzero integral first coordinate with $|u'|<u$ and preserves the
reduced denominator of the second. The first integrality follows from
$(Q-1)y\in\Z$; denominator preservation follows from
$\gcd(2Q-1,m)=1$. The norm gives
$u<y<(2Q-1)u/[2(Q-1)]$, proving $0<u'<u$ after reflection.
Descent reaches $u=y=1$, forcing $m=1$. Therefore
\[
 QV^2-(Q-1)y^2=1,\qquad z=\varepsilon gy,
 \quad\varepsilon\in\{-1,1\}.
\]
With $I=2Ug$, equations \eqref{eq:positive-trace} give
$f^2-\Delta I^2=1$, $Q=(\Delta I)^2$, $c^2\mid I$ and $c\le I$.
The retained signed all-index folding theorem applies to $V\equiv-c\pmod f$
and gives $c=\psi_A(p)$ for a positive odd index $p$. Since $c\ge3$,
$p\ge3$, and hence
\begin{equation}\label{eq:positive-c3}
 c\ge c_3:=4A^2-1=4\Delta+3>4q^4,
 \qquad c>X,\qquad K>0.
\end{equation}
No sign of $R,B,C,b,V$ or divisibility $q\mid Y$ was used.

\subsection{A uniform bound removes the positive shift}
If $b=0$, the auxiliary norm would give
$0<y^2-c^2=(c^2-1)/(Q-1)<1$, impossible for integers.
Thus $b\ne0$, $\sgn V=\sgn b$ and $|V|>1$.
Put $S=\Delta I$, so $Q=S^2$. Pell classification of $(SV,y)$
gives an odd auxiliary index: even traces are units modulo $S$,
whereas $SV$ is divisible by $S$. Index one would give $|V|=1$.
The index is therefore at least three, and $|V|\ge4Q-3$.
Moreover
\[
 \frac{a^2}{g^2}=\frac{\Delta U^2a}{1+a}>
 \frac{\Delta U^2}{2}>9,
\]
so $f>6g$. Using $c_3\le c$ and $f>2c$ yields the essential
bound, for either sign of $V$,
\begin{equation}\label{eq:positive-uniform}
 |b|\ge\frac{|V|-c}{f}
 \ge4\Delta f-\frac{c+c_3}{f}
 >4\Delta f-1>4g.
\end{equation}
The possible auxiliary index three is harmless; it is not excluded.
The norm also gives $y<3|V|/2$. Consequently
\[
 L=bV=|b||V|>0,\qquad gy<3L/8<L/2,
 \qquad |b|/q^d<L/2\quad(d=6,7).
\]
This remains valid at $Y=1$, before its cubic $q$-scale is recovered.

In the plain chart $e=k-R-1<b$ since $cT-k+1>0$.
In the inverse charts $k\ge2$ gives $e_*\le e<b$.
Thus $\rho<|b|/q^d$ in all charts. If $B\le-1$, then
$W_*t_*<-L$, and \eqref{eq:positive-carrier} implies
\[
 n=\varepsilon gy-\kappa-\sigma+W_*t_*+\rho
 <L/2-L+|b|/q^d<0,
\]
contradicting the positive domain. Hence $B\ge0$ and $C>0$.
This handles all negative values of $R$ and zero uniformly.
Only now use the authentic whole-tuple mask bound $3q<R<q^4-2$.
Equation~\eqref{eq:positive-c3} gives $c>2R$, hence $b,V>0$.
The retained normalized strong-rank and signed auxiliary-index theorems
apply to $V\equiv-c\pmod f$ and $V\equiv-R\pmod c$;
the latter uses $f-1=2\Delta i^2c^4E$. With $c>2p,2R$
they yield $p=R$. Thus $R$ is positive odd, and mask parity gives
even $q$ and odd $J$. Neither $\rho$ nor $\sigma$ has yet been
assumed integral.

\subsection{Authenticate the first index at every prime}
The square $E=\tau^2$ gives $k=2\psi_{1+2v}(m)$ with $m\ge1$.
The case $m=1$ would give $k=2$, $\eta=\zeta=1$ and
$c=2Y+1$ strictly between $\psi_A(1)=1$ and $\psi_A(2)=2A$.
Thus $m\ge2$. The first base $1+2v$ exceeds $A$, since
$\lambda Y(2X\lambda Y-X-1)>0$. Monotonicity and $c>k$
then give $m<R$.

In the free plain chart $v=q^4ws^2$. Integral $z$ and
$\gcd(q,hH)=1$ give $q^7\mid e-q^4s$, hence $q^4\mid e$.
The recurrence $k\equiv2m\pmod{4v}$ shows
$q^4\mid2m-R-1$. Its absolute value is less than $R<q^4$;
therefore $2m=R+1$.

In either inverse chart $v=q^3ws^2$. The inherited finite Lucas
identity, valid at every index, is
\[
 k=2m+8vC_m+32v^2Z_m,\quad
 C_m=\binom{m+1}{3},\quad Z_m\in\N.
\]
Since $k$ is even, $k(k^2-4)/6$ is integral: write $k=2u$
and use $3\mid u(u-1)(u+1)$. Thus $e_*$ is integral.
Clearing the carrier and cancelling its literal factor six, not
inverting six modulo $q$, gives $q^d\mid hH(e_*-q^4s)$,
hence $q^4\mid e_*$.
Put $t_0=2m$, $d_0=2m-R-1$, $A_m=C_m+4vZ_m$ and
$S_m=k^2+kt_0+t_0^2-4$. Cubic subtraction gives exactly
\begin{equation}\label{eq:inverse-cubic}
 3(e_*-d_0)=4v^2(24Z_m-A_mS_m).
\end{equation}
Indeed $t_0(t_0^2-4)=48C_m$ and $k-t_0=8vA_m$.
For any prime $p\mid q$, writing $a_p=v_p(q)\ge1$, the identity gives
\[
 v_p(e_*-d_0)\ge6a_p-v_p(3)\ge4a_p.
\]
This includes the prime three. Hence $q^4\mid d_0$, and the strict
size bound gives $d_0=0$. In all three charts,
\begin{equation}\label{eq:positive-index}
 m=(R+1)/2.
\end{equation}
Full $q^6$- or $q^7$-divisibility of $e_*$ has not been inferred here;
the third-primary obstruction is repaired only after projection.

\subsection{Projection precedes complete denominator recovery}
The literal numerator and integral $z$ give $H\mid hDK$.
Cancel only the $q$-part: $H\mid d_1hK$, where $d_1=1$ in
the plain chart and $d_1=6$ in the inverse charts.
Write $R=2r+1$, so $r\ge48$. The retained $\lambda$-uniform ratio
theorem applies to the authenticated indices and strip $Y<c/k<Y+1$:
\[
 Y>X^r/3,\qquad X^{r+1}<3Y(X+1)<H.
\]
Modulo odd $H$, the universal scaled recurrence gives
$3\cdot2^{R-1}\psi_A(R)\equiv4^R-1$. The retained projector yields
\[
 H\mid d_1h(X-2^R)(2^RX+1).
\]
The complete product, including $d_1$, has magnitude below $H$:
\[
 d_1h|(X-2^R)(2^RX+1)|
 <4\cdot4^rX^6+8\cdot16^rX^5<X^{r+1}<H.
\]
After division by $X^{r+1}$ the two terms decrease for $X\ge32$,
$r\ge48$; at $(32,48)$ their sum is $2^{-117}+2^{-25}<1$.
Thus $X=2^R$. The universal projector itself then gives positive
integral $\kappa=K/H$; no unproved cancellation of $d_1h$ is used.

Now $q\mid X$ implies $q=2^a$ with $a\ge5$.
Since $R>3q\ge6a$, one has $q^6\mid X$.
In the free charts $v=Xq^2Y^2$ is divisible by $q^8$.
In the plain chart $4v\mid e$ gives $q^7\mid e$; in the free
inverse chart \eqref{eq:inverse-cubic} with $d_0=0$ gives $q^7\mid e_*$.
In the scaled chart $Y=qs$, so $q^8\mid v$ as well and the same
identity gives $q^6\mid e_*$. Thus $\rho$ is integral everywhere.

The carrier now makes $\sigma-W_*B/h$ integral. Multiply by $q^3h$
and reduce modulo $q^3$ to obtain $q^3\mid Y$; reduce modulo $h$
to obtain $h\mid B$, using $\gcd(q,h)=1$ and $W_*\equiv1\pmod h$.
Thus $t_*=B/h\in\N$. These are the exact retained native outer
conditions. The fixed-minus parity and half-binomial theorems, uniform
in $\lambda$, reconstruct fresh source witnesses at
$A_{\rm source}=Y(X+1)+2$, not at the present base.
The retained whole-tuple native/refutation calibration then supplies a
finite syntactic ZFC refutation.

\subsection{A constructive positive lift with possibly negative correction}
Conversely a finite refutation supplies canonical outer data
$q=2^a$, $X=2^R$, $R>3q$, $R\equiv3\pmod4$, $q^3\mid Y$,
$C>0$ and integral $t_*=B/h\ge0$. Choose $w=X/q^2$ in the free
plain chart and $w=X/q$ otherwise; take $s=Y$ in free charts and
$s=Y/q$ in the scaled chart. All are positive integers.
Construct fresh $A,v$ and set
\[
 c=\psi_A(R),\quad k=2\psi_{1+2v}((R+1)/2),\quad
 \tau=\chi_{1+2v}((R+1)/2),
 \quad\eta=c-kY,\quad\zeta=k-\eta.
\]
The retained exact uniform ratio theorem gives $0<\eta<k$.
The same recurrence and power-of-two divisibility prove integral
$\rho$; projection gives $0<\kappa<Xc$, and $\sigma$ is positive
integral.

Choose a positive Pell matrix period at $A$ modulo $c^2\tau$,
obtaining a seed $(f_0,I_0)$ with $c^2\tau\mid I_0$.
Set $i=I_0/(c^2\tau)>0$. Then $a=f_0^2-1$, $G=(\tau f_0)^2$
and the doubled strong pair is $f=1+2a$, $I=2f_0I_0$.
Keep this seed fixed. The retained positive auxiliary reconstruction
gives $(V_*,y)$ with $V_*=f^2(cT_{\rm old}-R)-c$.
Define
\[
 T_{\rm new}=\frac{V_*+c+Rf}{cf}.
\]
This is integral: $f^2\mid V_*+c$, $V_*\equiv-R\pmod c$,
$f\equiv1\pmod c$ and $\gcd(c,f)=1$. It is strictly positive,
and $V_*=f(cT_{\rm new}-R)-c$ exactly. This is a fresh witness
construction, not an inverse shear on arbitrary positive tuples.

In inverse charts the correction may be negative. Since $k\ge R+1$
and $vk^2>6$, one has $|\rho|<vk^3/3$. For every chosen $i\ge1$,
\[
 g>\sqrt\Delta c^2E>c^2v^2k^2>vk^3,\qquad g>4Xc.
\]
For the latter inequality use $\sqrt\Delta>A-1>X$ and $cE>4$.
Thus $\kappa+\sigma<g/2$ and $|\rho|<g/3$.
In the plain chart $\rho\ge0$, and the same bound on $g$ suffices.
In all cases
\[
 n=gy+(1+hL)t_*+\rho-\kappa-\sigma
\]
is a strictly positive integer. Its carrier has $z=gy$, so the literal
polynomial vanishes. Every one of the ten independent inputs has been
constructed in the actual positive domain at one coherent tuple.

\subsection{Attained degrees and exact circuit counts}
Put $q_0=\beta J$, $k_0=\eta+\zeta$ and
$B_0=k_0sT-q_0^2(q_0-F)$. Complete homogeneous collection gives
\begin{align*}
 [P_{\rm plain}]_{240}
 &=-2^8q_0^{112}w^{26}s^{58}i^{10}k_0^{32}(q_0-F)^2,\\
 [P_{\rm scaled}]_{232}
 &=-9\cdot2^{10}q_0^{100}w^{26}s^{58}i^{10}k_0^{32}B_0^2,\\
 [P_{\rm free}]_{214}
 &=-9\cdot2^{10}q_0^{86}w^{26}s^{58}i^{10}k_0^{32}(q_0-F)^2.
\end{align*}
They have respectively $99,202,99$ collected monomials. In the free
charts the leading part of $b$ is $-q_0^3(q_0-F)$; in the scaled
chart the full degree tie $b_0=q_0B_0$ must be retained.
Selecting $\eta^{32}F^2$ in the free cases gives nonzero coefficients
$-2^8\beta^{112}$ and $-9\cdot2^{10}\beta^{86}$; selecting
$\eta^{34}T^2$ in the scaled case gives $-9\cdot2^{10}\beta^{100}$.
These prove attained degree at the authentic positive $\beta$.

In chart order, the degrees of $(Q,M,G,V,Z)$ are
$(82,13,52,40,60)$, $(80,11,50,40,57)$ and $(72,12,46,36,54)$.
The unique dominant output term is $-QM^2GV^2$.
The live binary-operation counts are respectively $82,85,87$.
The independent checker verifies complete top forms of every register,
not sampled lines. Historically, the scaled $232/85$ point saved two operations against
incoming Trace Radical's $232/87$; the free $214/87$ point saved eighteen
degrees against that same incoming point. Neither comparison is a global
bound or a claim about unrestricted integer variables.
